

%% packages utilises
%%---------------------
%\usepackage{etex}


\usepackage[utf8x]{inputenc}
\usepackage[T1]{fontenc}
\usepackage[english, french]{babel}
\usepackage[french]{translator}
\usepackage{frbib}
\usepackage{amsmath}
\usepackage{amssymb}
\usepackage{these}

\usepackage[a4paper]{geometry}

\usepackage{ulem}
%\normalem
%\usepackage{rotating}
\usepackage{tabularx}
\usepackage{textcase}
\usepackage{floatflt}
\usepackage{graphicx}
\usepackage{epstopdf}		% .eps to .pdf
\usepackage{moreverb} %% pour le verbatim en boite
\usepackage{multirow} %% pour regrouper un texte sur plusieurs lignes dans une table
\usepackage{url} %% pour citer les url par \url
\urlstyle{sf}
\usepackage[all]{xy} %% pour la barre au dessus des symboles
\usepackage{shorttoc} %% pour plusieurs tables des matières par la commande \shorttableofcontents{Titre}{profondeur}.
\usepackage{textcomp} %% pour le symbol pour mille par \textperthousand.

\usepackage[right]{eurosym}

\usepackage{latexsym}
% 
\usepackage{setspace} %allow to change line spacing
\usepackage{fancyhdr} %header / footer
  \setlength{\headheight}{36pt}
\usepackage[ Conny ]{ fncychap }
\usepackage{enumitem} %beautiful enumerations ... no beamer
\usepackage{subfig} %allow to split figure (article / book / report)
\usepackage{caption} %
\usepackage{color} %use color
\usepackage{epsfig,amsfonts} %math mode
\usepackage{array,colortbl} %use array / tabular and apply color / style on them
\usepackage{pstricks} %draw with latex
\usepackage{tikz}
\usepackage{listings} %source code with avec \being{lstlistgin}

\usepackage{stmaryrd}
%\usepackage{tocbibind}
%\usepackage[nottoc]{tocbibind}
\usepackage{pdfpages}
\usepackage{chngpage}
\usepackage{multicol}
\usepackage[fit]{truncate}
\pagestyle{fancy}
\fancyhead{} %on efface l'entête

\lhead{\includegraphics[width=1.2cm]{images/esi.jpeg}}
\chead{\scriptsize\textsc{Application des techniques de Machine Learning\\
à l’analyse de la conformité des transferts à l’étranger}}
\rhead{\includegraphics[width=1cm]{images/logo3.png}}
%\fancyhead[LE,RO]{\nouppercase{\truncate{0.5\headwidth}{\rightmark}}}
%\fancyhead[LO,RE]{\nouppercase{\truncate{0.5\headwidth}{\leftmark}}}
\fancyfoot{}
\renewcommand{\footrulewidth}{1pt}
\lfoot{\scriptsize\textbf{SEGHDA Ghislain}}
\rfoot{\scriptsize\textit\thepage}
\usepackage[Algorithme]{algorithm}
\usepackage{algpseudocode}

\usepackage{hyperref}
\usepackage[toc,nonumberlist]{glossaries}
\usepackage{wrapfig}

\usepackage[resetlabels]{multibib}
\newcites{publis}{Publications personnelles}

%% choix des profondeurs de section pour la table des matières
%% 2= subsection, 3=subsubsection
\setcounter{secnumdepth}{2}  %% Avec un numero.
\setcounter{tocdepth}{2}     %% Visibles dans la table des matieres

\makeglossary
\def\underscore{\char`\_}

\usepackage{nomencl}
%\renewcommand{\}{Liste des notations}
%\renewcommand*{\pagedeclaration}[1]{\unskip\dotfill\hyperpage{#1}}
\makenomenclature

\usepackage{makeidx}
\makeindex


\definecolor{colKeys}{rgb}{0,0,1}
\definecolor{colIdentifier}{rgb}{0,0,0}
\definecolor{colComments}{rgb}{0,0.5,1}
\definecolor{colString}{rgb}{0.6,0.1,0.1}

\definecolor{c1}{RGB}{219,144,71}
\definecolor{c2}{RGB}{100,212,78}
\definecolor{c3}{RGB}{255,111,111}
\definecolor{c4}{RGB}{111,139,255}

\definecolor{gris}{gray}{0.45}

\lstset{%configuration de listings
	float=hbp,%
	language=C,
	basicstyle=\ttfamily\small, %
	identifierstyle=\color{colIdentifier}, %
	keywordstyle=\color{colKeys}, %
	stringstyle=\color{colString}, %
	commentstyle=\color{colComments}, %
	columns=flexible, %
	tabsize=3, %
	%frame=trBL, %
	extendedchars=true, %
	%showspaces=false, %
	showstringspaces=false, %
	numbers=none, %
	breaklines=true, %
	breakautoindent=true, %
	captionpos=b,%
	xrightmargin=0cm, %
	xleftmargin=0cm,
	mathescape=true
}


\frenchspacing

\newgeometry{hmargin={0pt,0pt},vmargin={0pt,0pt}}
\savegeometry{include}

\newgeometry{vmargin={4.1cm,3.6cm},hmargin={3cm,2cm},twoside}
\setstretch{1.1}
\savegeometry{normalpreprint}

\newgeometry{vmargin={4.1cm,3.6cm},hmargin={2cm,3cm},twoside}
\setstretch{1.1}
\savegeometry{normalpreprintinverse}

\newgeometry{vmargin={4.1cm,3.6cm},hmargin={2.5cm,2.5cm}}
\setstretch{1.1}
\savegeometry{normaldigital}

\ifthenelse{\equal{\preprint}{YES}}
{
\loadgeometry{normalpreprint}
\savegeometry{normal}
\loadgeometry{normalpreprintinverse}
\savegeometry{normalinverse}
}
{
\loadgeometry{normaldigital}
\savegeometry{normal}
\savegeometry{normalinverse}
}


%%%%%%%%%%%%%%%%%%%%%%%%%%%%%%%%%%%%%%%%%%%%%%%%%%%%%%%%%%%%%
%	SubFloat environment for lstlisting inside subloat		%
%%%%%%%%%%%%%%%%%%%%%%%%%%%%%%%%%%%%%%%%%%%%%%%%%%%%%%%%%%%%%

\captionsetup[subfloat]{
}

\makeatletter
\newbox\sf@box
\newenvironment{SubFloat}[2][]%
	{	\def\sf@one{#1}%
		\def\sf@two{#2}%
		\setbox\sf@box\hbox
			\bgroup}%
	{	\egroup
	\ifx\@empty\sf@two\@empty\relax
		\def\sf@two{\@empty}
	\fi
	\ifx\@empty\sf@one\@empty\relax
		\subfloat[\sf@two]{\box\sf@box}%
	\else
		\subfloat[\sf@one][\sf@two]{\box\sf@box}%
	\fi}
\makeatother

\usepackage{bibentry}
\nobibliography*

%% macro/racourcis por les symboles et commandes usuelles
%%%%%%%%%%%%%%%% environnement pour les résumés%%%%%%%%%%%%%%%%%%%%
\makeatletter
\newenvironment{abstract}{%
    \null\vfil
    \@beginparpenalty\@lowpenalty
    \begin{center}%
            \@endparpenalty\@M
    \end{center}}%
   {\par\vfil\null}
\makeatother
% macro pour le 'debuggage': permet de corriger le document avec plus de facilit�
\newcommand{\mydebuglabel}[1]{\label{#1}}
%\newcommand{\debuglabel}[1]{\label{#1}}

% nouveaux environements de theoreme
%----------------------
\newtheorem{proposition}{Proposition}
%\newtheorem{lemme}{Lemme}
\newtheorem{exemple}{Exemple}
%\newcommand{\mydefinition}[2]{\begin{definition}{\bf #1:} #2$\diamond$\end{definition}}
%\newcommand{\mydefinition}[1]{\begin{definition}#1\end{definitio\newtheorem{theoreme}{Theorem}[section]
\newcommand{\deftitle}[1]{\begin{definition}{\bf (#1)}}

% alias de notation mathematiques
%-----------------------
%%%% debut macro %%%% %\newcommand{\barre}[1]{#1}  %% pour ajouter une barre horizontale au dessus d'un symbole

%\newcommand{\vector}[1]
%{\mathchoice
%{\overset{\mbox{\xymatrix{*{\hphantom{\displaystyle #1}}
%\ar[]+L;[]+R}}}{\displaystyle #1}}%
%{\overset{\mbox{\xymatrix{*{\hphantom{\textstyle #1}}
%\ar[]+L;[]+R}}}{\textstyle #1}}%
%{\overset{\mbox{\xymatrix{*{\hphantom{\scriptstyle #1}}
%\ar[]+L;[]+R}}}{\scriptstyle #1}}%
%{\overset{\mbox{\xymatrix{*{\hphantom{\scriptscriptstyle #1}}
%\ar[]+L;[]+R}}}{\scriptscriptstyle #1}}% 
%}

\newcommand{\barre}[1]{\overline{#1}}%% pour ajouter une barre horizontale au dessus d'un symbole
\newcommand{\prosite}[1]{{\bf{\small#1}}}%% pour �crire des motifs prosite
\newcommand{\permille}{\hbox{$\,^0\!/_{00}$}}%% symbole pourmille


%Pour changer la distance de la fl�che, on peut proc�der ainsi.
%\renewcommand{\ra}[1]
%{\overset{\raisebox{-1pt}{\mbox{\xymatrix{*{\hphantom{#1}}
%\ar[]+L;[]+R}}}}{#1}}
%%%% fin macro %%%%


%% alias specifiques pour les relations, utilisables seulement en mode mathematique
%-----------------------
\newcommand{\La}{\ensuremath{\langle}}
\newcommand{\Ra}{\ensuremath{\rangle}}

\newcommand{\dcp}[2]{\parallel_p^{\La #1,#2 \Ra}}
\newcommand{\dcs}[2]{\parallel_s^{\La #1,#2 \Ra}}
\newcommand{\cp}[2]{\parallel_p^{#1,#2}}
\newcommand{\cs}[2]{\parallel_s^{#1,#2}}
\newcommand{\cscc}{\parallel_s^{c_1,c_2}} % racourci pour cs avec comme classe c1 et c2
\newcommand{\scp}{\parallel_p} % scp pour Simple Common Prefix (simple = pas de classe)
\newcommand{\scs}{\parallel_s} % scs pour Simple Common Suffix (simple = pas de classe)
%\newcommand{\inc}[2]{\not\sim^{#1,#2}}
\newcommand{\sinc}{\parallel_s^{\neq}} % sinc pour Simple Incompatible (simple = pas de classe)
\newcommand{\stc}{\scs} % stc: suffixe toute classe
\newcommand{\scce}{\parallel_s^{ce}} % scce: suffixe commun de contre exemple
%\newcommand{\cp}[2]{{\vphantom{\parallel_p^{#2}}}^{#1}{\parallel_p^{#2}}}
%\newcommand{\cs}[2]{{\vphantom{\parallel_s^{#2}}}^{#1}{\parallel_s^{#2}}}

\newcommand{\suff}[2]{\text{{\it Suff}}_{#1}(#2)}
\newcommand{\pref}[2]{\text{{\it Pref}}_{#1}(#2)}
\newcommand{\relalias}{\ifmmode{\parallel_{r}}\else{$\parallel_{r}$}\fi} % un symbol pour repr�senter une relation qque entre 2 etats

%% alias pour la definition des classes d'automates, des acceptations et des automates particuliers (en mode math ou non)
%-----------------------
\newcommand{\assurepasmath}[1]{\ifmmode\text{#1}\else #1\fi}
%\newcommand{\cnfa}{\ifmmode\text{NFAC}\else NFAC}
\newcommand{\ufa}{\assurepasmath{UFA}}
\newcommand{\ufas}{\assurepasmath{UFAs}}
\newcommand{\dfa}{\assurepasmath{DFA}}
\newcommand{\dfas}{\assurepasmath{DFAs}}
\newcommand{\nfa}{\assurepasmath{NFA}}
\newcommand{\nfas}{\assurepasmath{NFAs}}
\newcommand{\rfsa}{\assurepasmath{RFSA}}
\newcommand{\rfsas}{\assurepasmath{RFSAs}}
\newcommand{\kamb}[1]{\ifmmode\text{AFA}_{#1}\else AFA\ensuremath{_{\text{#1}}}\fi}
\newcommand{\cnfa}{\assurepasmath{{NFC}}}
\newcommand{\cnfas}{\assurepasmath{{NFCs}}}
\newcommand{\cdfa}{\assurepasmath{{DFC}}}
\newcommand{\cdfas}{\assurepasmath{{DFCs}}}
\newcommand{\cufa}{\assurepasmath{{UFC}}}
\newcommand{\cufas}{\assurepasmath{{UFCs}}}
\newcommand{\ckamb}[1]{\ifmmode\text{{AFC}}_{#1}\else {AFC}\ensuremath{_{\text{#1}}}\fi}

\newcommand{\Sa}{\ifmmode\mathcal{S}\else\ensuremath{\mathcal{S}}\fi}

\newcommand{\auset}[2]{\ifmmode\boldsymbol{D}_{#1}(#2)\else\ensuremath{\boldsymbol{D}_{\text{#1}}(\text{#2})}\fi}
\newcommand{\partitions}[1]{\ifmmode\boldsymbol{P}(#1)\else\ensuremath{\boldsymbol{P}}(#1)\fi}
% ensemble des partitions restreinte a une classe d'automate
%\newcommand{\partitionsr}[2]{\ifmmode\mathcal{P}art_{#1}(#2)\else\ensuremath{\mathcal{P}art_{#1}}(#2)\fi}
\newcommand{\partitionsr}[2]{\ifmmode\boldsymbol{P}_{#1}(#2)\else\ensuremath{\boldsymbol{P}_{#1}}(#2)\fi}

\newcommand{\Acc}[2]{\ifmmode Acc_{#1}(#2)\else\ensuremath{Acc_{\text{#1}}(\text{#2})}\fi}
\newcommand{\smca}[1]{\ensuremath{\text{{\it MCA}}(}#1\ensuremath{)}}
\newcommand{\mca}{\ensuremath{\text{{\it MCA}}}} % mca sans le sample
\newcommand{\ua}{\ensuremath{\text{{\it UA}}}}
\newcommand{\sua}[1]{\ensuremath{\text{{\it UA}}(}#1\ensuremath{)}}
\newcommand{\mcak}[1]{\ensuremath{\text{{\it MCA}}_{#1}}}
\newcommand{\smcak}[2]{\ensuremath{\text{{\it MCA}}_{#1}(}#2\ensuremath{)}}
\newcommand{\smcau}[1]{\smcak{u}{#1}}
\newcommand{\mcau}{\mcak{u}}
\newcommand{\pta}{\ensuremath{\text{{\it PTA}}}}
\newcommand{\spta}[1]{\ensuremath{\text{{\it PTA}}(}#1\ensuremath{)}}
\newcommand{\au}[1]{\ensuremath{\La\Sigma,\Gamma_{#1},Q_{#1},I_{#1},\delta_{#1},\rho_{#1}\Ra}}% automate classifieur
\newcommand{\nau}[1]{\ensuremath{\La\Sigma,Q_{#1},I_{#1},\delta_{#1},F_{#1}\Ra}}% automate normal

%% alias des op�rateurs de parcours
%---------------------------------
% seulement en mode math !!!
% fusion simple sur partition
\newcommand{\ispar}{\prec} % inf�riorit� stricte
\newcommand{\ipar}{\prec^*} % inf�riorit�
\newcommand{\fuspar}{\xrightarrow{\prec}} % op�rateur de fusion
%\newcommand{\fispar}{\xrightarrow{\succ}}% op�rateur de fission
% fusion simple sur automates
\newcommand{\isaut}{\prec_{A}} % inf�riorit� stricte
\newcommand{\iaut}{\prec^*_{A}} % inf�riorit�
\newcommand{\fusaut}{\xrightarrow{\prec_A}}% op�rateur de fusion simple
\newcommand{\fisaut}{\xrightarrow{\succ_A}}% op�rateur de fission simple
\newcommand{\fussaut}{\xrightarrow{A*}}% op�rateur de suite de fusion simples
% fusion deterministe
\newcommand{\idet}{\prec^*_{det}}
\newcommand{\isdet}{\prec_{det}}
\newcommand{\fdet}{\xrightarrow{det}}
% fusion restreinte au determinisme
\newcommand{\irdet}{\prec^*_{D}}
\newcommand{\isrdet}{\prec_{D}}
\newcommand{\frdet}{\xrightarrow{D}}
% fusion desambiguisante
\newcommand{\ides}{\prec^*_{des}}
\newcommand{\isdes}{\prec_{des}}
\newcommand{\fdes}{\xrightarrow{des}}
% fusion restreinte a la k-ambiguite
\newcommand{\ikam}[1]{\prec^*_{#1-amb}}
\newcommand{\iskam}[1]{\prec_{#1-amb}}
\newcommand{\fkam}[1]{\xrightarrow{#1-amb}}

%% alias de presentation
%-----------------------
\newcommand{\etc}{etc.}
%\newcommand{\commentaire}[1]{{\tiny #1}}
\newcommand{\commentaire}[1]{}
\newcommand{\proof}[1]{{\bf Preuve:} #1}
%\newcommand{\proof}[1]{{\bf Proof:} #1$\square$}
\newcommand{\hintproof}[1]{{\bf Id�e de la preuve:} #1}
%\newcommand{\hintproof}[1]{{\bf Hint of the proof:} #1$\square$}
%\newcommand{\mycaption}[1]{\caption{{\small #1}}}
%\newcommand{\mycaption}[1]{\caption{ #1}}

%\newcommand{\remark}[1]{{\small Remark: \emph{#1}}}
\newcommand{\remark}[1]{Remark: \emph{#1}}

\newcommand{\resumefr}[1]{%
\thispagestyle{empty}
\section*{R�sum�}
#1 
\vfill
}

% >> macro pour la separation d'un bout de page en deux parties, utile pour les figures et leur caption a droite ou gauche
% l'argument specifie la taille de la partie gauche
\newlength\jataille
\newcommand{\figgauche}[3]%
{\jataille=\textwidth\advance\jataille by -#1
\parbox{#1}{#2}
\parbox{\jataille}{#3}
}

\newcommand{\figtxt}[1]%
{{\it {\small #1}}}
% << fin macro

% >> macro pour tracer un trai horizontal sur la largeur de la page
\newcommand{\traithoriz}{\raisebox{0.4em}{\vrule depth 0pt height 0.4pt width \textwidth}}


% nom des algos
\newcommand{\edsm}{EDSM}
\newcommand{\rpni}{RPNI}
\newcommand{\edsmdfaf}{D$_{\text{f}}$\textit{edsm}}
\newcommand{\edsmdfac}{D$_{\text{c}}$\textit{edsm}}
\newcommand{\hcdfaf}{D$_{\text{f}}$\textit{hc}}
\newcommand{\hcdfac}{D$_{\text{c}}$\textit{hc}}
\newcommand{\hcufaf}{U$_{\text{f}}$\textit{hc}}
\newcommand{\hcufac}{U$_{\text{c}}$\textit{hc}}
\newcommand{\delete}{DLT2}
\newcommand{\majvote}{MAJ}


\newcommand{\kcite}[1]{
\begin{verse}
\cite{#1}~\bibentry{#1}
\end{verse}
}
%\newcommand{\kcite}[1]{\cite{#1}}

%%%% debut macro pour faire des lignes �paisses dans les tableaux %%%%
\makeatletter
\def\hlinewd#1{%
\noalign{\ifnum0=`}\fi\hrule \@height #1 %
\futurelet\reserved@a\@xhline}
\makeatother
%%%% fin macro %%%%


%%%%%%%%%%%%%%%%%%%%%%%%%%%%%%%%%%%%%%%%%%%%%%%%%%%%%%%%%%
%%%%%%%%%%%%%%%%%		Exemple		%%%%%%%%%%%%%%%%%%%%%%
%%%%%%%%%%%%%%%%%%%%%%%%%%%%%%%%%%%%%%%%%%%%%%%%%%%%%%%%%%

% glossary entry : defines the identifier in the glossary table (1rst arg), 
% and the description in second arg.
\newglossaryentry{id}{
name={ID}, %what appear in text
description={\emph{Identifier} : Identifiant}, %what appear in glossary
first={ID (\emph{Identifier} : Identifiant)}} %what appear the first time in text
\newcommand{\id}{\gls{id}} % shortcut pour the glossary entry. Simply coll \id in the text
\newcommand{\ID}{ID} % to put the text without using the glossary shortcut.


\newglossaryentry{sg}{
name={SG}, %what appear in text
description={Société Générale}, %what appear in glossary
first={SG ( Société Générale)}} %what appear the first time in text
\newcommand{\sg}{\gls{sg}} % shortcut pour the glossary entry. Simply coll \id in the text

\newglossaryentry{sgbf}{
name={SGBF}, %what appear in text
description={Société Générale Burkina Faso}, %what appear in glossary
first={SGBF ( Société Générale Burkina Faso)}} %what appear the first time in text
\newcommand{\sgbf}{\gls{sgbf}} % shortcut pour the glossary entry. Simply coll \id in the text




%%  1ere de Couverture (titre sur une seule ligne latex)
\titre{Mémoire de fin de cycle master SI-SAD}

\soutenue
%%   Laisser cette ligne en commentaire sauf pour la version finale.
%%   (la premiere page contiendra "a soutenir le ..."
%%   au lieu de "soutenue le ...")


%% Les différents champs de la couverture...
\datesout{`` fin 20xx ''}
\Auteur{Hoda}{DABONNE}%sur une seule ligne latex

\ecole{Ecole Supérieur d'Informatique}  % { rennes1 | insa | ens }
\Labo{Lab.}
\Umr{XXXX}
\LaboEtendu{Laboratoire de recherche} % sauf INSA
\ComposanteUniversitaire{Université Nazi Boni} % sauf INSA

%% La composition du jury : prénom, nom, titre
\President[Mme]{Prenom}{Nom}{Titre} %% le président du jury
\Rapporteur{Prenom}{Nom}{Titre}
\Rapporteur[Mme]{Prenom}{Nom}{Titre}
%% Si vous avez N rapporteurs ça marche toujours...
\Examinateur[Mme]{Prenom}{Nom}{Titre}
\Examinateur{Prenom}{Nom}{Titre}
%% idem...
\Advisor{Serge}{Miguet}{Maitre de Conférences}
\Advisor{Prénom}{Nom}{Titre}

\Mention{Thèmatique}
\Keywords{Mots clés}
\ordre{00000}    %% le numéro d'ordre donné par la Scol.


\title{\TITRE}

\author{\AUTp \AUTn}

\hypersetup{
	pdfauthor={{\AUTp} {\AUTn}},
	pdftitle={{\TITRE}},
	pdfsubject={\MENTION},
	pdfkeywords={\KEYWORDS},
	hidelinks
}
